\sectionframe{Setups and tasks}{What we predict, and how we check it}

\begin{frame}{The same data supports many tasks}
  \begin{center}
    \renewcommand{\arraystretch}{1.15}
    \small
    \begin{tabular}{llll}
      \toprule
      \head{Task} & \head{Input} & \head{Output} & \head{Example} \\
      \midrule
      \alert{Forecasting} & the past of a series & its future & bikes tomorrow \\
      \alert{Classification} & a whole series & a label & sleep stage of an EEG epoch \\
      Regression & a whole series & a number & age from an ECG \\
      Anomaly detection & a series & flagged times & a broken counter \\
      Imputation & a series with gaps & the missing values & offline weeks \\
      Clustering & many series & groups & commute or leisure counters \\
      \bottomrule
    \end{tabular}
  \end{center}
  \takeaway{Part 1 covers forecasting, parts 2 and 3 cover classification.}
\end{frame}

\begin{frame}{Forecasting predicts the future from the past, up to a horizon}
  \begin{columns}[c]
    \column{.6\textwidth}
    \figtodo{one forecast on Sébastopol, context and horizon shaded (notebook \S2)}
    \column{.36\textwidth}
    \begin{itemize}
      \item \textbf{Context}: the history the model reads.
      \item \textbf{Origin}: the last observed time.
      \item \textbf{Horizon}: how many steps ahead.
      \item \textbf{Lead time}: the gap from the origin to a given step.
    \end{itemize}
  \end{columns}
\end{frame}

\begin{frame}{A point forecast hides the uncertainty, so we predict distributions}
  \begin{columns}[c]
    \column{.6\textwidth}
    \figtodo{fan chart with the 10, 50 and 90\% quantiles (notebook \S2)}
    \column{.36\textwidth}
    \begin{itemize}
      \item The median as the point forecast.
      \item Quantiles give prediction intervals.
      \item Sample paths give joint scenarios.
    \end{itemize}
  \end{columns}
  \takeaway{A forecast without its uncertainty cannot support a decision.}
\end{frame}

\begin{frame}{Forecasting setups multiply: univariate, multivariate, with covariates, across series}
  \begin{columns}[T]
    \column{.48\textwidth}
    \begin{block}{Univariate}
      One counter, its own past only. \muted{Notebook \S4}
    \end{block}
    \begin{block}{Multivariate}
      Both directions of one street, forecast jointly. \muted{\S7}
    \end{block}
    \column{.48\textwidth}
    \begin{alertblock}{With covariates}
      Calendar known ahead, weather known only in the past. \muted{\S8, \S9}
    \end{alertblock}
    \begin{alertblock}{Many series}
      The 128 counters of Paris at once. \muted{\S10}
    \end{alertblock}
  \end{columns}
\end{frame}

\begin{frame}{Local models fit each series, global models learn across series}
  \begin{columns}[T]
    \column{.48\textwidth}
    \begin{exampleblock}{Local}
      ARIMA, ETS: one small model per series, fitted on its own history.
    \end{exampleblock}
    \column{.48\textwidth}
    \begin{block}{Global}
      One model trained on a panel of series, shared across all of them.
    \end{block}
  \end{columns}
  \vskip2mm
  A foundation model is a global model trained on a very wide panel, then used on
  series it never saw.
  \takeaway{The wider the panel, the less each new series has to teach.}
\end{frame}

\begin{frame}{Classification assigns a label to a whole series}
  \begin{itemize}
    \item Activity from the accelerometer of a phone.
    \item Sleep stage from a 30 s window of EEG.
    \item Type of a variable star from its light curve.
  \end{itemize}
  \vskip2mm
  The model reads the whole window at once, there is no future to predict.
  \takeaway{Parts 2 (Mantis) and 3 (EEG) are about this task.}
\end{frame}

\begin{frame}{Evaluation must respect time}
  \begin{columns}[c]
    \column{.6\textwidth}
    \figtodo{the 2025 forecast origins of the notebook on a timeline (\S4)}
    \column{.36\textwidth}
    \begin{itemize}
      \item A random split leaks the future into training.
      \item Rolling origin: forecast from several dates, with only the past
        before each one.
    \end{itemize}
  \end{columns}
  \takeaway{Test on the future, from many origins.}
\end{frame}

\begin{frame}{Scaled metrics make series comparable}
  \begin{itemize}
    \item \textbf{MASE}: the error relative to the seasonal naive forecast,
      $\displaystyle \frac{\frac1H\sum_h |y_{T+h}-\hat y_{T+h}|}{\frac1{T-m}\sum_{t>m} |y_t - y_{t-m}|}$.
      Below 1 beats the baseline.
    \item \textbf{WQL}: the quantile loss summed over the quantile levels $q$,
      $\rho_q(y, \hat y_q) = \max\big(q(y-\hat y_q), (q-1)(y-\hat y_q)\big)$,
      scaled by $\sum |y|$. It rewards sharp and calibrated quantiles.
    \item \textbf{Coverage}: how often the truth falls in the 80\% interval.
      It should be close to 0.8.
  \end{itemize}
  \takeaway{Point accuracy and calibration are measured separately.}
\end{frame}

\begin{frame}{Always compare with the seasonal naive baseline}
  \begin{columns}[c]
    \column{.5\textwidth}
    \begin{center}
      \begin{tabular}{lc}
        \toprule
        \head{Forecast, 2-week horizon} & \head{MASE} \\
        \midrule
        Seasonal naive (last week again) & 1.155 \\
        Best foundation model & \alert{0.926} \\
        \bottomrule
      \end{tabular}
    \end{center}
    \column{.44\textwidth}
    Boulevard de Sébastopol, hourly, rolling origins in 2025.
    \muted{Notebook \S4}
  \end{columns}
  \takeaway{A gain of 20\%, not an order of magnitude.}
\end{frame}
